\documentclass[12pt]{article}
\usepackage[utf8]{inputenc}
\usepackage[T1]{fontenc}
\usepackage{hyperref}
\title{Mbonga Language Resources for K-12 Curriculum Development}
\author{Compiled by Claude Sonnet 5.5 (claude-sonnet-5-5) for Andr\'as Kornai}
\date{}
\begin{document}
\maketitle

\section{Introduction}
Mbonga (also called Mboa; ISO 639-3 xmb; Glottocode mbon1252) is a Niger-Congo (Southern Bantoid) language of the East region of Cameroon, west of B\'etar\'e-Oya. Sources disagree on its status. Wikipedia calls it ``apparently extinct'' and cites 1,490 speakers (2000). The Endangered Languages catalogue rates it threatened, with low certainty. Joshua Project treats the Mbonga as a living people group of about 3,100 in Cameroon. The searches found no school materials, no writing system and no audio recordings. This report is therefore short.

\section{Primary Educational Resources}
No primers, textbooks or school programmes for Mbonga were found in the search results. The Endangered Languages catalogue reports no writing system.

\section{Linguistic and Reference Materials}
\begin{itemize}
\item Wikipedia, Mboa language (status, classification, speaker figure): \url{https://en.wikipedia.org/wiki/Mboa_language}
\item Endangered Languages catalogue entry (threatened, no writing system): \url{https://endangeredlanguages.com/lang/6031}
\item Numerals dataset (Max Planck Institute, Mbonga page; notes a possible vigesimal system and that data above ten is needed): \url{https://lingweb.eva.mpg.de/channumerals/Mbonga.htm}
\end{itemize}

\section{Religious and Community Resources}
\begin{itemize}
\item Joshua Project, Mbonga in Cameroon (population and Bible translation status): \url{https://www.joshuaproject.net/index.php/people_groups/13630/CM}
\end{itemize}

\section{Digital and Audio Resources}
\begin{itemize}
\item Global Recordings Network language page for Mbonga, which lists no recordings: \url{https://globalrecordings.net/en/language/xmb}
\end{itemize}

\section{Curriculum Development Context}
There are almost no documented resources, so any heritage or revitalization curriculum would first need fieldwork and community-led documentation. Speaker status should be verified with the community, because the sources conflict. The catalogue entries above are the only leads found.
\end{document}